
\begin{table*}[!th]
    \centering
    \small
    \caption{
        \textbf{Prompt used to check \theorem-Theorem documents.} {\color{blue} Blue text} denote instance-specific inputs. 
        For each question, we provide the problem statement, the solution, the theorem name, the theorem definition, and the document text from a ProofWiki document.
    }
    \label{tab:theorem_check}
    % \resizebox{0.98\linewidth}{!}{
    \begin{tabular}{>{\raggedright\arraybackslash\tt}p{0.98\textwidth}<{}}
        \toprule
            \vspace{-1em}
            Instruction: Determine if the given text can be helpful in understanding and solving the given problem. Use the theorem, its definition, and the problem solution to help you make a decision. The text can be helpful if it uses very similar reasoning steps as the solution and applies the theorem in a related way as the solution applies the theorem to solve the problem. The text should be able to assist a student who is not familiar with the theorem in ultimately solving the problem.\\\\

            The input is given to you in a json format with the following keys:\\
            Problem statement: The problem statement that the student is trying to solve.\\
            Solution: The solution to the problem.\\
            Theorem: The theorem that is used in the solution.\\
            Theorem definition: The definition of the theorem.\\
            Text: The text that you need to evaluate.\\

            Think step by step and reason about the theorem and the text first before finally making a decision. Output True if the text is helpful, and False if it is not.
            \\
            {\color{blue}\{
                "Problem statement": "",
                "Solution": "",
                "Theorem": "",
                "Theorem definition": "",
                "Text": "",
            \}}
            \\
            Now, write your answer in the following format:\\
            Reasoning: [your reasoning here]\\
            Answer: [True/False]\\

       \bottomrule
    \end{tabular}
\end{table*}
